% TOP_PROBLEM: {"id": 30000859, "title": "RSW Theory for Anisotropic Square-Lattice Percolation", "category_id": 19, "category": {"id": 19, "name": "probability", "display_name": "Probability", "description": "Problems involving probability theory, stochastic processes, and random structures.", "slug": "probability", "order_index": 19, "created_at": "2026-07-31T15:26:25.671Z"}, "status": "open"}
% FIRST_POSTED: 2026-10-06
% ATTEMPT_STATUS: unresolved
% OPENAI_CLAIM: {"schema_version":1,"source_repo":"https://github.com/openai/math","source_commit":"adc7f1241b42e322a6451854ab7e4b4c146bf78a","release_date":"2026-10-06","match":"partial","resolution":"partial","independently_reviewed":false,"second_pass":"n/a","families":[{"family":"223","title":"Random-cluster interfaces: critical, disordered, thermal, and natural-time scaling","subject":"Probability and statistical mechanics","manuscripts":[{"title":"Conformal Limits of Critical Square-Lattice Random-Cluster Interfaces","pdf_path":"preprints/Conformal-Limits-of-Critical-Square-Lattice-Random-Cluster-Interfaces-October-5-2026/paper.pdf","readme_path":"preprints/Conformal-Limits-of-Critical-Square-Lattice-Random-Cluster-Interfaces-October-5-2026/README.md","date":"2026-10-05","theorem_ref":"intro:main (sections/introduction.tex:125)","statement_sources":[{"path":"preprints/Conformal-Limits-of-Critical-Square-Lattice-Random-Cluster-Interfaces-October-5-2026/build/sections/introduction.tex","line":125,"label":"intro:main"}]}],"lean":null,"reasoning_trace":null}]}
\documentclass[11pt]{article}
\usepackage[margin=25mm]{geometry}
\usepackage{fontspec}
\setmainfont{Latin Modern Roman}
\usepackage{amsmath,amssymb,mathtools}
\usepackage{unicode-math}
\setmathfont{Latin Modern Math}
\usepackage{enumitem}
\usepackage{xurl,hyperref,bookmark}
\hypersetup{colorlinks=true,urlcolor=blue}
\setcounter{secnumdepth}{0}
\setlength{\parindent}{0pt}
\setlength{\parskip}{5pt}
\setlength{\emergencystretch}{3em}
\begin{document}
\section*{RSW Theory for Anisotropic Square-Lattice Percolation}

\subsection{OpenAI's claim}
OpenAI presents a mathematical result relevant to this question in the selected manuscript. The later interface theorem extends the square-lattice FK claim through $q=4$, with the same SLE parameter relation and a separate complete nested-loop limit.

\subsection{Mathematical statement of the claimed result}
\paragraph{Result 1.} Fix \(q\in[1,4]\), \(p_q=\sqrt q/(1+\sqrt q)\), and
\[
 \kappa(q)=\frac{4\pi}{\arccos(-\sqrt q/2)}.
\]
Use the square-lattice random-cluster law with edge-subset weights \(p_q^{|A|}(1-p_q)^{|E|-|A|}q^{k_\xi(A)}\), where \(\xi\) wires one boundary arc and leaves the other free. For a bounded Jordan domain with two marked boundary points, approximate its boundary uniformly by simple lattice polygons, preserving the marks. The claim gives
\[
 \eta_n\ \Longrightarrow\ \operatorname{SLE}_{\kappa(q)}
\]
for the Dobrushin interface in uniform curve distance modulo increasing reparametrization. Under the free infinite-volume law, its full nested loop collection satisfies
\[
 \Gamma_\delta\ \Longrightarrow\ \text{nested whole-plane }\operatorname{CLE}_{\kappa(q)}
\]
in the spherical matching topology: all loops above a shrinking diameter threshold are matched with shrinking uniform spherical error. Both limits hold along every mesh sequence, with \(q\) fixed.

\subsection{Outline of the argument}
The tiled comparison and angular-transfer method is supplemented by endpoint estimates at $q=4$. Boundary current tests identify the driving process, while full exploration and localization retain macroscopic loops; at $q=1$ this also yields the bond-percolation crossing formula.

\subsection{Scope relative to this problem}
At q=1 the source covers isotropic critical square-lattice bond percolation (p\_h=p\_v=1/2), yielding crossing control in this restricted subcase. It does not cover general anisotropic p\_h+p\_v=1 or determine the requested general shape class. This entry records partial support; complete coverage of the question is not established by this catalogue.

\subsection{Formal verification}
Family 223 has no Lean scope page in this release. Its selected result is not presented as formalized. This site has not executed the Lean code. The scope descriptions do not certify the remaining manuscript arguments.

\subsection{Status caveat}
OpenAI attributes these manuscripts to work by an internal model. Its release cautions that arguments without formalization can contain mistakes. This catalogue has not independently checked the proofs; the labels report the mathematical claims attributed to OpenAI.

\subsection{Sources}
\begin{itemize}
\item OpenAI, \emph{Conformal Limits of Critical Square-Lattice Random-Cluster Interfaces}, 2026-10-05; citation key \texttt{OAI:Conformal-Limits-of-Critical-Square-Lattice-Random-Cluster-Interfaces-October-5-2026}. \href{https://github.com/openai/math/blob/adc7f1241b42e322a6451854ab7e4b4c146bf78a/preprints/Conformal-Limits-of-Critical-Square-Lattice-Random-Cluster-Interfaces-October-5-2026/paper.pdf}{PDF (GitHub)}; \href{https://github.com/openai/math/raw/adc7f1241b42e322a6451854ab7e4b4c146bf78a/preprints/Conformal-Limits-of-Critical-Square-Lattice-Random-Cluster-Interfaces-October-5-2026/paper.pdf}{PDF (direct)}; \href{https://github.com/openai/math/blob/adc7f1241b42e322a6451854ab7e4b4c146bf78a/preprints/Conformal-Limits-of-Critical-Square-Lattice-Random-Cluster-Interfaces-October-5-2026/README.md}{Manuscript page and citation}; \href{https://github.com/openai/math/blob/main/preprints/Conformal-Limits-of-Critical-Square-Lattice-Random-Cluster-Interfaces-October-5-2026/paper.pdf}{Latest version}; Statement source: \href{https://github.com/openai/math/blob/adc7f1241b42e322a6451854ab7e4b4c146bf78a/preprints/Conformal-Limits-of-Critical-Square-Lattice-Random-Cluster-Interfaces-October-5-2026/build/sections/introduction.tex\#L125}{intro:main}.
\item \href{https://github.com/openai/math/blob/adc7f1241b42e322a6451854ab7e4b4c146bf78a/CONTENTS.md}{OpenAI catalogue entry}.
\end{itemize}
\end{document}
